\documentclass[11pt,letterpaper]{article}
\usepackage[margin=0.85in]{geometry}
\usepackage{amsmath,amssymb,booktabs,array,longtable,graphicx}
\usepackage{xcolor,enumitem,fancyhdr,hyperref}
\hypersetup{colorlinks=true,linkcolor=blue!45!black,urlcolor=blue!45!black,
 citecolor=blue!45!black,pdftitle={Ember: September 10 working paper},
 pdfauthor={Ember project, prepared for Gregory Laughlin}}
\definecolor{ember}{HTML}{A4492A}
\setlength{\parindent}{0pt}
\setlength{\parskip}{6pt}
\setlength{\emergencystretch}{2em}
\setlist{nosep,leftmargin=*}
\pagestyle{fancy}\fancyhf{}
\fancyhead[L]{\small\textcolor{ember}{EMBER}\quad Working paper}
\fancyhead[R]{\small 10 September 2026}
\fancyfoot[C]{\small\thepage}\setlength{\headheight}{14pt}
\newcommand{\Ms}{M_\odot}\newcommand{\Rs}{R_\odot}\newcommand{\Ls}{L_\odot}
\newcommand{\Tyr}{\mathrm{Tyr}}
\newcommand{\doc}[2]{\href{https://github.com/oklo/ember/blob/master/#1}{#2}}
\begin{document}
\begin{center}
{\LARGE\bfseries\textcolor{ember}{Ember}}\\[6pt]
{\Large Low-mass stellar evolution toward\\an ultra-cold helium remnant}\\[8pt]
{\normalsize Working paper and reproducibility checkpoint}\\[3pt]
{\small 10 September 2026; completed stellar results through $3.40\Tyr$}
\end{center}

\textbf{Current result.} The fixed-baryonic-mass $0.1\Ms$ reference has reached
$3.40\times10^{12}$ yr from its specified static initial model. It remains
fully convective, with $X_{\rm H}=0.2033564$ and $T_{\rm eff}=3201.95$ K.
Hydrogen exhaustion, a helium white dwarf, and remnant cooling have not yet
been obtained. The explicit final objective is a supported, evolved cooling
sequence through $T_{\rm eff}=1000$, 500, and 100 K. A fitted cooling law or
an extrapolated atmosphere would not meet that objective. We use
$1\Tyr=10^{12}$ yr throughout.

Since the September 9 edition, the selected hydrogen-poor EOS and radiative
opacity families have been refined and independently checked; a larger gas
atmosphere has enabled the new track; plasma-neutrino losses have entered the
energy equation; and exact nuclear caches have reduced measured computing
cost. Independent F77 reproductions now locate their own radiative-core
transitions. Preparatory cold-electron and dense-ion audits address parts of
the remnant problem, but do not constitute an installed remnant EOS.

The immediate atmospheric limit is now $X_{\rm H}=0.2$, rather than 0.3.
A candidate extending to 0.1 has complete source and runtime support checks,
but an independent comparison exposes composition-interpolation curvature.
Refinement and a separate stellar diagnostic run are in progress; they are
not included as completed results in this edition.

\section{Reference model and completed trajectory}
The initial composition is $X_{\rm H}=0.7$, $X_3=0$, $Z=0.02$, with the GS98
elemental metal pattern. Age excludes formation and pre-main-sequence
contraction. Rotation, magnetic fields, external heating, and mass loss are
absent from this reference. The precision shown below identifies saved
states; it is not a statement of comparable physical accuracy.

\begin{center}\begin{tabular}{@{}lrr@{}}\toprule
Quantity & $3.30\Tyr$ & $3.40\Tyr$ \\\midrule
Radius ($\Rs$) & 0.14695441 & 0.14588637 \\
Photon luminosity ($\Ls$) & 0.0019776154 & 0.0020155006 \\
Effective temperature (K) & 3175.197 & 3201.953 \\
Central temperature (MK) & 7.767698 & 8.031406 \\
Central density ($\mathrm{g\,cm^{-3}}$) & 252.2519 & 258.2795 \\
Central hydrogen fraction & 0.22229873 & 0.20335640 \\
Central $^3$He fraction & 0.001357904 & 0.000970265 \\
Remaining hydrogen mass ($\Ms$) & 0.02222987 & 0.02033564 \\
Convective mass fraction & 1 & 1 \\\bottomrule
\end{tabular}\end{center}

A fresh 512-point calculation reached $3.30\Tyr$ with 1822 accepted steps
and three rejections. Its exact native restart added 38 accepted steps to
reach $3.40\Tyr$, with no additional rejection. The joined history contains
1861 distinct states. The copied executable and physical inputs remained
fixed; source-tree edits during the calculation did not alter that executable.
The previous $2.85\Tyr$ gas reference and its recovery inputs are preserved.

\begin{figure}[ht]\centering
\includegraphics[width=\textwidth]{evolution_history.pdf}
\caption{Completed 0--$3.40\Tyr$ trajectory using the selected 96-node gas
atmosphere, refined 64-plane EOS and TOPS v4 opacity. The fresh 0--3.30T
history is joined to its exact 3.30--3.40T restart after checking identical
overlapping physical states. Dotted lines mark the selected atmospheric
hydrogen floor and temperature ceiling. The current $^3$He depletion and
radius decrease occur during continued hydrogen burning.}
\end{figure}

\section{Implemented physics and numerical scope}
Ember is a C++23 development of the stellar-evolution line reconstructed in
\href{https://github.com/oklo/Henyey}{Henyey/F77}. The evolutionary picture of
Laughlin, Bodenheimer \& Adams \cite{lba97} is a benchmark to investigate;
its lifetime is not imposed on the new calculation. The two codes use
different physical prescriptions and numerical methods. Greater computational
cost does not by itself establish greater physical accuracy.

\begin{longtable}{@{}>{\raggedright\arraybackslash}p{0.19\textwidth}>{\raggedright\arraybackslash}p{0.75\textwidth}@{}}
\toprule Component & Current treatment and limits \\\midrule\endhead
EOS & FreeEOS 3.0 EOS1 represented by a common Helmholtz potential, with
consistent caloric and transport responses and explicit H/$^3$He composition
dependence. The selected GS98 family has 64 planes, including true zero H
within its original temperature/density domain. This is not cold dense-matter
coverage. Potassium is absent in the source; its omitted mass fraction is
$4.56\times10^{-6}$. Helium isotope and partition-function approximations remain. \\[5pt]
Opacity & AESOPUS 2.1 low-temperature gas opacity and LANL TOPS hot/dense
opacity, with supported blends, explicit source-family selection, and isotope
number-density mapping. Refined hydrogen-poor TOPS v4 is selected. Substituted
source cells and unsupported density/composition stencils are rejected. \\[5pt]
Burning & Non-equilibrium $^3$He in a reduced ppI/ppII network; Solar Fusion II
S-factor quadrature and finite-degeneracy Salpeter--Van Horn screening.
Atomic mass defects determine deposited heat and nuclear-neutrino energy.
Pep, hep, ppIII and a CNO network remain absent; metals are inert. \\[5pt]
Transport & MLT with Schwarzschild or full-EOS Ledoux buoyancy; approximate
mixtures of Ioffe/Cassisi conduction inputs. The reference selects Ledoux,
semiconvective and thermohaline composition mixing. Microscopic diffusion,
settling, and a separate semiconvective heat flux are not implemented. \\[5pt]
Atmosphere & Offline TLUSTY208/SYNSPEC54 LTE plane-parallel gas atmospheres,
matched at $\tau=100$. Molecular/atomic absorption, CIA and scattering enter
the source models. Runtime interpolation uses H, $^3$He, $T_{\rm eff}$ and
gravity, with explicit masks for missing nodes. No accepted condensate family. \\[5pt]
Thermal losses & The selected HRW plasma-decay prescription \cite{hrw94}
enters central and shell energy equations with analytic derivatives. Pair,
photo, bremsstrahlung, recombination and partial-ionization corrections are
not provided by this selector. \\[5pt]
Evolution & Analytic structure Jacobians; damped Henyey/Newton solves;
coupled implicit burning and conservative mixing; step-doubling time control.
The mass mesh is fixed. Native restarts bind internal state and controller
state to executable bytes, input tables, tolerances and physical selections. \\
\bottomrule\end{longtable}

The material EOS, atmosphere chemistry, opacities and conduction do not all
derive from one thermochemical source. Nuclear screening is not derived from
the FreeEOS interaction potential. Source agreement and mathematical derivative
consistency therefore require separate assessment from physical accuracy.

\section{Input refinement and remaining atmospheric limits}
\subsection{Selected interior inputs}
The first zero-H EOS extension failed its heat-capacity acceptance with a
2.367\% discrepancy in cool trace-hydrogen molecular gas. Additional composition
planes reduced curvature error. The accepted 64-plane family passed 528 fresh
source comparisons: maximum relative differences in $P$, $E$, $c_V$ and $c_P$
are respectively 0.05213\%, 0.10875\%, 0.18413\%, and 0.19216\%.
All 2048 sampled legacy profile queries retain the old results exactly.
See \doc{docs/results/metal_eos_exhaustion_refined_v2_audit.json}{the EOS audit}.

The coarse TOPS interpolation had an 11.90\% discrepancy at an independent
source midpoint. Successive separately identified refinements have reduced
the largest discrepancy in the latest independent, defined active-domain
source-node audit to 1.7062\%. AESOPUS is held fixed. This number is neither
a stellar-profile error nor a global opacity uncertainty between source nodes.
The original failed comparison is retained. The selected family and audit
definitions are recorded in
\doc{docs/results/opacity_hydrogen_poor_refined_v4_audit.json}{the TOPS v4 audit}.

\subsection{Selected 96-node grid and candidate 108-node extension}
The 96-node selected atmosphere adds hydrogen-poor and warmer source states
to the prior 72-node gas family. Its global axes are
$X_{\rm H}=.2,.3,.45,.7$, $X_3=0,.12$,
$T_{\rm eff}=2600,2800,3000,3200,3400$ K, and
$\log g=4.9,5.15,5.4$. These axes do not imply a filled rectangle:
missing source states are masked, and runtime evaluation requires complete
interpolation and derivative stencils.

Twelve completed $X_{\rm H}=.1$ source states at $T_{\rm eff}=3200,3400$ K
produce a 108-node candidate. All 108 runtime source-node queries and their
EOS inversions are supported; 112 retained warm-grid queries preserve values
and derivatives exactly. Source/interior density differences range from
$-0.01082\%$ to $+1.99104\%$ across candidate nodes.

Independent lower-boundary controls test actual achieved optical depths,
not nominal requested depths. At the two $X_{\rm H}=.1$, 3200 K,
$\log g=5.4$ corners, depths approximately 165/339 and 170/350 give maximum
matching-state fractional changes $7.13\times10^{-6}$ and $1.93\times10^{-5}$.
TLUSTY's TAULAS does not truncate a supplied mass grid; verified seed-column
truncation is therefore explicit and independently checked.

\begin{center}\begin{tabular}{@{}lrr@{}}\toprule
Independent $X_{\rm H}=.15$, $X_3=.005$, 3300 K, $\log g=5.1$
 & $\Delta T/T$ & $\Delta P_{\rm gas}/P_{\rm gas}$ \\\midrule
Full 108-node interpolation & $+2.13408\%$ & $-0.78677\%$ \\
Composition interpolation at fixed $T_{\rm eff},g$ & $+1.65199\%$ & $-0.78424\%$ \\
Temperature/gravity contribution at composition corners & $+0.47426\%$ & $-0.00254\%$ \\\bottomrule
\end{tabular}\end{center}

Four newly completed canonical source controls separate these contributions.
They add in logarithms, with closure below $9\times10^{-16}$; fractional
differences multiply. This diagnostic identifies hydrogen spacing as the
dominant matching-temperature discrepancy. New $X_{\rm H}=.15$ node planes
and fresh $X_{\rm H}=.125,.175$ heldouts are being calculated. The original
$X_{\rm H}=.15,X_3=.005$ source remains an independent point, not a promoted
node. At the older independent $X_{\rm H}=.25$ point the differences remain
$+0.69036\%$ in matching temperature and $-0.93414\%$ in pressure.

The 108-node candidate is used only by a separately identified diagnostic
stellar run toward 3.60T while this refinement proceeds. Its result is not
included in the completed 3.40T reference. Warm gas source convergence does
not validate a condensate boundary or dense cold helium atmosphere.

\subsection{Condensates, line broadening and zero hydrogen}
The condensate production guard remains active. Previously converged solar
2800 K, low-gravity source corners failed the grain-enthalpy check; an
artificial extra heat capacity is a sensitivity test, not complete chemistry.
A new equilibrium correction uses neutral-atomic chemical potentials to
evaluate the difference in $g/T$ between gas-only and gas-plus-condensate
equilibria. Derivatives provide volume, enthalpy, entropy and heat capacity,
including rearrangement of gas molecules and ions. At one audited state the
complete $c_P$ correction is $117496\ \mathrm{erg\,g^{-1}\,K^{-1}}$, compared
with the earlier grain-only estimate of 289000. No correction is installed;
bulk normalization, a common baseline and phase-boundary responses remain
necessary. A corrected calling-convention audit found no history dependence
in either the old or defensive-copy adapter's twelve fixed-bulk calls.

The warmer material-opacity ceiling is about 10103 K. A 15000 K pilot
encountered nonfinite extrapolation of the native He I 4471 profile. A reduced
bottom column is useful only when its full profile has source support and
matching-state depth convergence; it cannot justify extrapolating line data.

Zero H is supported in the interior composition table but not yet in the
atmosphere generator. Its abundance normalization and the native molecular
initialization depend on hydrogen. A direct 4017 K opacity-source diagnostic now confirms that relabelling
all number abundances by a common factor is insufficient: electron densities
agree to about $1.2\times10^{-6}$ relative, but individual sampled absorption
coefficients change substantially (up to a factor 27.3 for scale 0.15).
This single-isotherm result does not measure a stellar-boundary error; it
rejects an unverified normalization change. Relabelling tiny nonzero H as
zero would not establish the requested endpoint.

\section{Convection and the F77 comparison}
At 3.40T the minimum offline nodal ratio of diffusive (radiative plus
conductive) to adiabatic gradient is 1.33593, at mass fraction 0.211. At 3.30T
it was 1.58798. The central ratios are 3.08980 and 3.70705 respectively.
Composition is homogeneous, so unity is the local stability threshold.
These nodal measurements suggest approach to loss of full convection;
they do not demonstrate central radiative-core formation or determine an
onset age. The runtime uses its own face-based mixing criterion.

Separate read-only-source F77 reproductions add a print-only core diagnostic.
An unmodified source control reproduces the original outputs byte for byte
for the full AJR run and a 6000-model Ferguson control. The comparison uses
151 initial mesh points, $\alpha=1.36$, explicit input decks and the retained
source; it is not a new measurement of the original 1997 paper's model.

\begin{center}\begin{tabular}{@{}lrr@{}}\toprule
F77 opacity selection & First central nonconvective flag (Tyr) & $X_{\rm H,c}$ \\\midrule
AJR & 3.22353 & $1.75623\times10^{-2}$ \\
Ferguson & 2.94629 & $1.06427\times10^{-5}$ \\\bottomrule
\end{tabular}\end{center}
F77's convective flag has hysteresis; the AJR gradient crossed unity one
accepted model earlier, at 3.22314T. The Ferguson evolution later lost
convergence during cooling at about 2.95471T. Its earlier accepted onset
bracket remains valid, but its cooling run is incomplete.

At matched hydrogen abundance near 0.203, AJR F77 is about 74.9\% more
luminous than Ember: $L\simeq0.003525\Ls$ versus $0.0020155\Ls$.
The surface temperatures are about 3698 and 3202 K, while radii and central
temperatures agree within 0.5\%. The corresponding ages are approximately
2.50737T and 3.40T; comparing models at equal age would mix different
compositional stages. F77's $^3$He abundance is about 56\% higher.

Changing only the F77 opacity selector to Ferguson increases luminosity by
54.5\% at this H abundance. At fixed central temperature, density and
H/$^3$He coordinates, the two pp-heating prescriptions differ by roughly
9\%, rather than the whole 75\% surface-luminosity gap. These are evidence
for atmosphere/opacity and the resulting composition history as leading
contributors. They do not uniquely separate EOS, screening, boundary and
coupling effects. The grey-boundary tendency discussed by Baraffe et al.
\cite{baraffe98} supports the direction of a possible boundary effect, not
its magnitude here. F77 already includes CNO, conduction and adaptive-mesh
capabilities; Ember is not a strict superset of its physics.

\section{Energy accounting and cold-matter preparation}
The selected discrete energy balance is
\[
 L_\gamma=L_{\rm nuc,dep}+L_{\rm grav}-L_{\nu,\rm thermal},\qquad
 L_{\rm nuc,rest}=L_{\rm nuc,dep}+L_{\nu,\rm nuclear}.
\]
At 3.40T these luminosities, in solar units, are
$L_{\rm nuc,dep}=0.0020154351$,
$L_{\rm grav}=6.5481\times10^{-8}$,
$L_{\nu,\rm nuclear}=4.08865\times10^{-5}$ and
$L_{\nu,\rm plasma}=3.98645\times10^{-11}$.
Gravothermal luminosity refers to the last accepted implicit half step;
nuclear and thermal-neutrino diagnostics use the reported state.
The discrete luminosity-balance residual is below $10^{-15}$ at this endpoint.
That checks the implemented discrete equations, not an independent exact
global gravitational-energy integral or the absence of missing physics.

Cold-electron thermal derivatives originally suffered cancellation between
large degenerate contributions. Paired quadrature on opposite sides of the
Fermi surface now retains the small response. Pure-He tests at 100--10000 K
and $\rho=10^3$--$10^6\ \mathrm{g\,cm^{-3}}$ recover the Sommerfeld limit;
at 100 K the heat capacity agrees to approximately $10^{-10}$ relative.
The largest tested deviation from the leading limit is $8.20\times10^{-7}$
at 10000 K, where finite-temperature terms remain resolved. The current
source also supplies electron occupation entropy and ionic translational
entropy including mixing, tested against the first law and Maxwell relation.
These analytic components are separate from the production FreeEOS potential.

An isolated Ioffe EOSFI22 audit compares explicitly forced liquid and solid
free energies and their derivatives at exact density. It found and locally
repaired two quantum pressure-derivative defects; original source and patch
are retained. Across sixteen forced-phase states, the largest density-response
defect decreases from 6.77 to $1.05\times10^{-6}$ in ideal-ion pressure units
after finite-difference refinement. This validates internal derivatives of
the formulae, not the stable phase or their physical accuracy.

Helium quantum melting, mixture corrections, latent heat, caloric reference
matching and joins to partially ionized matter remain unfinished.
The source's default classical $\Gamma=175$ switch must not be mistaken for
a demonstrated helium phase calculation \cite{baiko22}. No dense remnant EOS
has been installed in the stellar evolution. The detailed scope is in
\doc{docs/DENSE_EOS.md}{DENSE\_EOS.md} and
\doc{docs/NEUTRINOS.md}{NEUTRINOS.md}.

\section{Acceptance criteria for the remaining evolution}
\begin{enumerate}[itemsep=5pt]
\item \textbf{Composition and atmospheric coverage.} Finish the diagnosed
warm-grid refinement with independent heldouts, then extend along measured
stellar profiles toward zero H and higher temperature/gravity. Resolve
condensate enthalpy and unsupported line broadening without clipping.
\item \textbf{Stratification and burning.} Add conservative microscopic
diffusion and settling, including the ambipolar electric field and diffusion
energetics. Resolve moving convective boundaries and thin fuel layers;
introduce conservative remeshing where fixed-mesh convergence fails.
Assess semiconvective heat transport, pep/ppIII/CNO contributions and
screening/EOS consistency under the newly encountered conditions.
\item \textbf{Distinct exhaustion endpoints.} Retain brackets and recrossings
of central $X_{\rm H}=10^{-3},10^{-4},10^{-5}$, and separately of deposited
burning luminosity below 10\%, 1\% and 0.1\% of the photon luminosity on the
cooling branch. Record residual hydrogen mass and shell burning. A core
crossing alone does not establish extinction of all hydrogen burning.
\item \textbf{Dense matter and cooling boundary.} Couple liquid/solid helium
thermodynamics and mixture physics to the interior; validate heat capacities,
coexistence and conductivity in the actual regime. Follow surface H/He from
transport. Dense cool atmospheres require nonideal chemistry, pressure-dependent
absorption and CIA, with condensate, refraction and dense-fluid relevance
checks \cite{blouin18}. Ordinary low-density gas opacity is insufficient.
\item \textbf{Convergence and uncertainty.} Converge exhaustion and 1000/500/100 K
cooling endpoints in mesh, timestep and source-grid resolution. Preserve
fresh-start and exact-restart controls. Separate interpolation error,
discretization error and uncertainty in the physical prescriptions. Assess
baryonic versus gravitating mass before quoting precision lifetimes.
\end{enumerate}

The wider programme includes mass and initial-composition surveys with an
explicit $Y(Z)$ relation and metal pattern, wind/rotation/magnetic
sensitivities, and distinct minimum-mass and higher-mass evolutionary branches.
Helium enrichment by burning at fixed $Z$ is not a survey of stars born at
high $Z$. No lifetime maximum or historical nominal duration is prescribed.
Environmental heating, encounters and hypothetical particle-physics effects
belong to separately stated future scenarios \cite{adams97}.

\section{Computing cost, reproducibility and publication}
Atmosphere generation is reusable source work: expensive frequency-dependent
radiative transfer and molecular chemistry are computed offline and tabulated.
Further work is required when a track leaves validated composition, temperature
or gravity coverage, or when source physics/resolution changes. The stellar
integrator interpolates those tables; it does not solve every wavelength at
every timestep.

The 0--3.30T star used 2350.57 child CPU seconds and 4088.19 elapsed UTC
seconds with competing source jobs. The exact 3.30--3.40T continuation used
97.00 CPU seconds. CPU and elapsed clocks are reported separately because
concurrent jobs and system suspension affect wall time. Two isolated ABBA
benchmarks of exact nuclear caches reduced CPU by 12.52\% and 10.84\%, with
byte-identical complete outputs; their product is about 22\% for those
benchmark conditions, not a whole-lifetime performance claim.

The separate unmodified 151-point F77 AJR reproduction used 27.37 CPU seconds
through its native cooling stop. Ember's 512-point coupled implicit solves,
step doubling and different physical prescriptions perform additional work.
No matched-input full-star profile has uniquely apportioned the performance
gap. An older atmosphere replay placed about 54\% of its phase time in
chemistry and just 0.0033 seconds in the dense linear solve; that is one
source replay, not a general stellar bottleneck measurement. Details and
receipts are in \doc{docs/COMPUTATIONAL_COST.md}{COMPUTATIONAL\_COST.md}.

GitHub carries source, generators, configurations, small provenance reports,
the compact history CSV and this paper. Newly generated bulk tables,
raw histories, native checkpoints and archived executables remain local.
The adjacent \texttt{recovery\_manifest.json} records the raw/compressed
hashes of the 3.30T and 3.40T histories and receipts, latest native checkpoint
and exact copied executable. Its ignored \texttt{artifacts/} directory is
not bundled in a fresh clone. September 9 recovery artifacts are retained
locally; the September 10 edition replaces the old paper in the current tree.

The history figure rebuilds from the committed CSV alone. With local archives,
\texttt{build\_figures.py --from-archives} rechecks the seam and rebuilds the
CSV. The adjacent \texttt{validation.json} records build, test and document
checks. Table-dependent tests require the locally installed inputs described
in \doc{docs/DATA_REPRODUCTION.md}{DATA\_REPRODUCTION.md}.

Native restart requires identical executable and physical-input bytes,
tolerances and mesh. New physics therefore starts a separately identified
calculation; old checkpoints are not silently reinterpreted. The 3.40T
reference executable has SHA-256
\begin{quote}\small\ttfamily
24a6ec3de0d06bf1e58e4ff4527fa4b38f4b02774be1703cba4ba371d7f15237.
\end{quote}
The latest live-job state is recorded at the beginning of
\doc{HANDOFF.md}{HANDOFF.md}; dated reports and saved running markers alone
cannot certify process liveness. Inspect existing processes and completed
receipts before launching further work.

\clearpage
\begin{thebibliography}{99}\small
\bibitem{lba97} Laughlin, G., Bodenheimer, P., \& Adams, F. C. (1997).
The End of the Main Sequence. \emph{ApJ}, 482, 420.
\href{https://doi.org/10.1086/304125}{doi:10.1086/304125}.
\bibitem{adams97} Adams, F. C., \& Laughlin, G. (1997).
A dying universe: the long-term fate and evolution of astrophysical objects.
\emph{Rev. Mod. Phys.}, 69, 337.
\href{https://doi.org/10.1103/RevModPhys.69.337}{doi:10.1103/RevModPhys.69.337}.
\bibitem{baraffe98} Baraffe, I., Chabrier, G., Allard, F., \& Hauschildt, P. H. (1998).
Evolutionary models for solar metallicity low-mass stars: mass--magnitude
relationships and color--magnitude diagrams. \emph{A\&A}, 337, 403.
\href{https://arxiv.org/abs/astro-ph/9805009}{arXiv:astro-ph/9805009}.
\bibitem{hrw94} Haft, M., Raffelt, G., \& Weiss, A. (1994).
Standard and nonstandard plasma neutrino emission revisited.
\href{https://arxiv.org/abs/astro-ph/9309014}{arXiv:astro-ph/9309014}.
\bibitem{baiko22} Baiko, D. A., \& Chugunov, A. I. (2022).
Ab initio thermodynamics of one-component plasma for astrophysics of white
dwarfs and neutron stars. \emph{MNRAS}, 510, 2628.
\href{https://doi.org/10.1093/mnras/stab3613}{doi:10.1093/mnras/stab3613}.
\bibitem{blouin18} Blouin, S., Dufour, P., \& Allard, N. F. (2018).
A New Generation of Cool White Dwarf Atmosphere Models. I. Theoretical Framework
and Applications to DZ Stars. \emph{ApJ}, 863, 184.
\href{https://arxiv.org/abs/1807.06616}{arXiv:1807.06616}.
\end{thebibliography}
\end{document}
